\documentclass[../../main.tex]{subfiles}
\begin{document}

\begin{center}
  {\fontsize{48pt}{58pt}\selectfont \textbf{2003}\normalsize\textbf{年度}}
\end{center}

\begin{table}[htb]
  \centering
  \caption{出題テーマと難易度}
  \begin{tabular}{|c|p{6.2cm}|c|c|c|}
    \hline
    \textbf{問題番号} & \textbf{テーマ} & \textbf{思考量} & \textbf{計算量} & \textbf{総合} \\
    \hline
    第1問 & 正三角形上の点と重心 & 易 & 易 & 易 \\
    \hline
    第2問 & 正三角形を折ってできる四面体の最大体積 & 普 & 普 & 普 \\
    \hline
    第3問 & 複素数解の距離が等しい三次方程式 & 難 & 難 & 難 \\
    \hline
    第4問 & 正の数列の不等式 & 易 & 易 & 易 \\
    \hline
    第5問 & 交代和の極限 & 普 & 普 & 普 \\
    \hline
    第6問 & 図形に従って文字を並べる場合の数 & 易 & 易 & 易 \\
    \hline
  \end{tabular}
\end{table}

\begin{table}[htb]
  \centering
  \caption{解答}
  \renewcommand{\arraystretch}{1.4}
  \begin{tabular}{|c|c|p{9cm}|}
    \hline
    \textbf{問題番号} & \textbf{小問} & \textbf{答え} \\
    \hline
    第1問 & -- & 証明問題 \\
    \hline
    第2問 & -- & $D$は$AC$の中点，最大体積$\dfrac{\sqrt{3}}{48}$ \\
    \hline
    第3問 & -- & $(a,b)=(0,0),\left(3\sqrt[3]{2},3\sqrt[3]{4}\right)$ \\
    \hline
    第4問 & -- & 証明問題 \\
    \hline
    第5問 & -- & $\dfrac{1}{2}$ \\
    \hline
    第6問 & -- & $6\cdot3^6=4374$通り \\
    \hline
  \end{tabular}
\end{table}

\paragraph{試験形態}
6問．試験時間は要確認．

\end{document}
